\section{Introduction}
\label{sec:introduction}

A rigid frictional system can lose forward uniqueness at a Painlev\'e
singularity. Compliance restores an ordinary evolution, but its
vanishing limit need not select a single outgoing motion. We study
a compliant, smooth-friction unfolding of the critical oscillator
of Nordmark, V\'arkonyi and Champneys~\cite[Section~6.1]{NVC2018}.
Two masses are joined by a rigid link, with the lower mass contacting
a moving belt. Successive contacts and flights form a cascade whose
number of returns tends to infinity as compliance vanishes.

Our question concerns two different forms of independence from the
auxiliary contact parameters. Do they give the same family of limiting
paths? Can one sequence of compliances select the same member of that
family for all parameter values? We answer both questions for paths
stopped at the first crossing of a fixed nearby link angle. The system
has relaxing normal memory and friction smoothed at slip scale $\eps$,
with belt speed $U_b=U_*-k\sqrt\eps$ and fixed $k>0$.
Section~\ref{sec:model} specifies the compact mechanical, profile and
physical-preparation classes used throughout.

\paragraph{Earlier work and the present problem.}
Champneys and V\'arkonyi~\cite[Sections~4, 5.1 and~5.3]{CV2016}
place contact regularization, dynamic jam and reverse chatter within
the forward-continuation problem for frictional rigid bodies.
Compliance-dependent continuation already appears in reverse
chatter: Nordmark, Dankowicz and
Champneys~\cite[Sections~5.2 and~6.1]{NDC2009} give an asymptotic
matching argument for selecting different forward motions along
different compliance sequences, with numerical terminal curves for
a conservative normal law and Coulomb friction. The relaxing
normal law used here was introduced by Nordmark, V\'arkonyi and
Champneys~\cite[Equations~(20)--(21)]{NVC2018} in their study of
dynamic jam. The concluding discussion of their 2017 preprint
identifies motion beyond the first lift-off or the onset of impact
without collision as a further problem, including repeated
contacts. We determine that subsequent local motion under the
smooth friction and belt scaling fixed below, including every
stopped path and the common-sequence question.

Rigorous passage results in related contact problems resolve
several of the local difficulties. Hogan and
Kristiansen~\cite[Theorems~3.1--3.2]{HK2017} prove compression,
slip arrest and lift-off for a single impact without collision,
and a threshold separating it from immediate lift-off, at a fixed
negative Painlev\'e coefficient. Kristiansen and
Hogan~\cite[Theorem~6]{KH2018} prove persistence of a strong
canard at the degenerate point, using a second blow-up to resolve
the ellipticity left by the first. Their
erratum~\cite{KH2018Erratum} requires sufficiently small initial
tangential velocity for rigid-model slip arrest before the singular
point; the canard theorem is unchanged. The additional tasks here
are preparation of the independent normal memory into a repeatable
return domain, comparison through a growing number of returns,
and actual realization of every stopped contact prefix.

Regularization can also select one outgoing orbit in other
nonsmooth systems. For a truncated visible--invisible two-fold
normal form, Kristiansen and Hogan~\cite[Theorem~1]{KH2019TwoFold}
prove uniform contraction of a transverse incoming family towards
one outgoing orbit under arctangent regularization, with a
nonresonance condition and separation from the weak canard;
their Section~5 discusses other regularization functions.
In three-dimensional contact mechanics, Cheesman, Hogan and
Kristiansen~\cite[Sections~4--5]{CHK2023} use normal compliance
and a separate smoothing of spatial Coulomb friction to analyse
impact without collision. They resolve local sticking dynamics
and identify a Painlev\'e-I equation in the transition from
slipping to inconsistency.

For homogeneous singular ODEs, Drivas and
Mailybaev~\cite[Theorem~5 and Corollary~2]{DM2018} study expelling
regularizations approaching an exponentially attracting
renormalized cycle with positive mean radial drift.
Regularizations escaping into the basin of the same cycle share
a phase family of limiting continuations. Their geometric sequence
$\nu_n=\exp\{-T\langle F_r\rangle n+\chi\}$ selects phase
$\zeta=c-\chi$, with regularization-dependent $c$.
This gives a related distinction between a common limiting family
and a common selecting sequence. In the contact problem here, an
entrance-phase difference disappears at a free target but remains
larger than the fixed-angle contact stopping window. Selection
therefore depends on the target type, and the order of the entrance
constants determines the endpoint limits.

The scalar phase of a power-like return has a classical precedent.
Goncharuk and Kudryashov~\cite[Lemmas~5--6]{GK2019} prove
rectification and refined double-logarithmic parameter asymptotics,
following root estimates of Ilyashenko, Kudryashov and
Schurov~\cite[Sections~4.5--4.6]{IKS2018} and the smooth saddle-loop
work of Dumortier and Roussarie~\cite{DR1990}. We construct the
phase together with the complete mechanical graph and control its
accuracy along actual returns. Bressan and
Mazzola~\cite[Theorem~2]{BM2013} prove time--state graph convergence
for specified approximations of impulsive feedbacks. Here the
contact curves and their ordering are derived from the compliant
force law.

\paragraph{Main conclusions.}
Theorem~\ref{thm:whole} classifies and realizes every stopped-path
limit from the prescribed initial data. Each path consists of its
incoming rigid arc, an infinite ordered contact--flight cascade
leaving the singular point, and a terminal contact or flight prefix.
At fixed mechanical parameters and observation section, the outgoing
family is common to the prescribed auxiliary classes; different
incoming data retain their own incoming arcs.
Theorem~\ref{thm:common-free} gives the stronger common-sequence
conclusion for every free-flight target, uniformly over those compact
classes.

Theorem~\ref{thm:contact-selection} states the entrance-constant law
for contact targets within the same physical model and classes.
Two fixed unequal constants exclude common interior contact selection,
even at different contact phases. Moreover, on any sequence selecting
an interior contact target for one datum, smaller and larger entrance
constants force admission and release for the other data.
Proposition~\ref{sel:threshold-selection} gives explicit threshold
sequences and uniform endpoint selection on compact subclasses
separated from the threshold. Equal constants are necessary for
common interior selection; these estimates do not make them
sufficient. The belt application is
Corollary~\ref{cor:belt-contact-selection};
Corollary~\ref{cor:friction-profile-selection} gives the corresponding
application along an admissible friction-profile curve constructed
locally through each baseline profile.

\paragraph{Mechanism.}
Recurring compression consumes tangential slip; arrest changes the
normal response, and the inherited compression drives release.
The subsequent cubic free-flight balance gives a return scale
of order $\delta^{3/4}$. On the invariant graph of the complete
normalized limiting return, the exact base map and its phase satisfy
\[
 F_\lambda(\delta)\sim K_\lambda\delta^{3/4},\qquad
 \mathcal S_\lambda(F_\lambda(\delta))
       =\tfrac34\mathcal S_\lambda(\delta).
\]
The prescribed physical preparation determines a bounded entrance
correction. Repeated return contraction reduces the difference of
two such corrections to order $1/\log(1/\eps)$ at a fixed positive
section. An actual contact stop at $p=-as$ requires the much narrower
window $\delta=s+O_s(\eps)$. Thus
\[
 \eps\ll\frac1{\log(1/\eps)}\ll1.
\]
The same entrance difference vanishes in the free-target limit but
still separates the contact stopping windows and orders their
endpoint limits. The comparison uses the exact phase at the fixed
section, not just its leading logarithmic approximation.

\paragraph{Proof route.}
Section~\ref{sec:local-results} derives the complete prepared entrance
and the repeatable local return from the physical initial data.
Section~\ref{sec:graph} constructs the exact graph and compares the
actual full states through the required event indices.
Section~\ref{sec:physical-paths} pairs whole returns and scans a
fixed actual contact index to classify and realize complete paths.
Section~\ref{sec:contact-selection} compares the residual phase with
the contact window and proves the belt and profile applications;
Section~\ref{sec:observations} treats longer paths, first-angle
truncation and strict local crossings. The local estimates and
cancellations remain in the appendices.
